\documentclass[a4paper, 12pt]{extarticle}

% Поддержка языков
\usepackage{polyglossia}
\setdefaultlanguage[forceheadingpunctuation=false]{russian}
\setotherlanguage{english}

% Настройка шрифтов
\usepackage{fontspec}
\usepackage{mathtools}
\usepackage{unicode-math}
\setmainfont[Ligatures=TeX]{NewCM10-Book} % Шрифт для основного текста документа
\setsansfont[Ligatures=TeX]{NewCMSans10-Book}
\setmonofont{NewCMMono10-Book} % Шрифт для кода
\setmathfont{NewCMMath-Book}

% Настройка отступов от краев страницы
\usepackage[left=3cm, right=1.5cm, top=2cm, bottom=2cm]{geometry}

\usepackage{titleps} % Колонтитулы
\usepackage{subfig} % Для подписей к рисункам и таблицам
\usepackage{graphicx} % для вставки картинок
\graphicspath{{./img/}} % Путь до папки с изображениями
\usepackage{wrapfig2}
\usepackage{tikz} % Пакет для отрисовки графиков
\usetikzlibrary{arrows,positioning,shadows}
\usepackage{stmaryrd} % Стрелки в формулах
\usepackage{indentfirst} % Красная строка после заголовка
\usepackage{hhline} % Улучшенные горизонтальные линии в таблицах
\usepackage{multirow} % Ячейки в несколько строчек в таблицах
\usepackage{longtable} % Многостраничные таблицы
\usepackage{paralist,array} % Список внутри таблицы
\usepackage{lua-ul} % Зачеркнутый, подчеркнутый текст
\usepackage{amsmath, amsthm} % ams пакеты для математики, табуляции
\usepackage{nicematrix} % Особые матрицы pNiceArray
\usepackage{icomma} % Убрать пробел после запятой как десятичного разделителя

\usepackage{esvect} %Адаптивные векторы

\usepackage{float}

\usepackage{caption}

\captionsetup[figure]{labelsep=period}  % Точка после номера рисунка
\captionsetup[table]{labelsep=period}   % точка после номера таблицы

% Стили теорем и формул, нумерация вида (2.1), (3.5)
\numberwithin{equation}{section}
\numberwithin{figure}{section}
\numberwithin{table}{section}

\theoremstyle{plain}
\newtheorem{theorem}{Теорема}[section]
\newtheorem{lemma}[theorem]{Лемма}
\newtheorem{corollary}{Следствие}[theorem]

\theoremstyle{definition}
\newtheorem{definition}{Определение}[section]

\theoremstyle{remark}
\newtheorem{remark}{Замечание}[section]

\theoremstyle{definition}
\newtheorem{problem}{Задача}[section]

\theoremstyle{remark}
\newtheorem*{solution}{Решение}

\usepackage{enumitem} % Расстояние между элементами списка
\setlist{nosep}

% работа с русскими обозначениями элементов списка \begin{enumerate}[label=\asbuk*), ref=\asbuk*]
\makeatletter
\AddEnumerateCounter{\asbuk}{\russian@asbuk@alph}{щ}
\makeatother

\usepackage[singlespacing]{setspace} % Межстрочный интервал
\setlength{\parindent}{1.25cm} % Табуляция
\setlength{\parskip}{0cm}


% Добавляем гипертекстовое оглавление в PDF
\usepackage[
bookmarks=true, colorlinks=true, unicode=true,
urlcolor=blue,linkcolor=black, anchorcolor=black,
citecolor=black, menucolor=black, filecolor=black,
]{hyperref}

\usepackage{xurl}

% Настройка переносов
\usepackage[protrusion=true, expansion=true]{microtype} % Для LuaLaTeX microtype использует font expansion, что заметно уменьшает число переносов и делает край более ровным.

\usepackage{cleveref}

\newpagestyle{main}{
	% Верхний колонтитул
	\setheadrule{0cm} % Размер линии отделяющей колонтитул от страницы
	\sethead{}{}{} % Содержание {слева}{по центру}{справа}
	% Нижний колонтитул
	\setfootrule{0cm} % Размер линии отделяющей колонтитул от страницы
	\setfoot{}{\thepage}{} % Содержание {слева}{по центру}{справа}
}
\pagestyle{main}

% НОВЫЕ КОМАНДЫ
\newcommand{\n}{\par}
\newcommand{\dd}{\mathrm{d}}
\newcommand{\pp}{\partial}
\newcommand{\M}{\mathrm{M}}
\newcommand{\Rey}{\mathrm{Re}}

\DeclareMathOperator{\rot}{rot}
\DeclareMathOperator{\diverg}{div}
\DeclareMathOperator{\ctg}{ctg}
\DeclareMathOperator{\grad}{grad}
\DeclareMathOperator{\tg}{tg}

% Правильная норма с автоматическим размером $\norm*{\dfrac{f(x)}{g(x)}}$
\DeclarePairedDelimiter\norm{\lVert}{\rVert}

% --- Компенсация артефактов OTF-математики у интегралов с пределами ---
\newcommand*{\intgap}{\mskip-9mu\relax} % стягивание материала после оператора
\newcommand*{\intlowshift}{8mu}         % сдвиг нижнего предела влево
\newcommand*{\intupshift}{3mu}          % сдвиг верхнего предела влево

\NewDocumentCommand{\opwithlimits}{m t\limits e{_^}}{%
	#1%
	\IfBooleanTF{#2}{% ветка с \limits: пределы сверху/снизу + все правки
		\limits
		\IfValueT{#3}{_{\mkern-\intlowshift\relax#3\mkern\intlowshift\relax}}%
		\IfValueT{#4}{^{\mkern-\intupshift\relax#4\mkern\intupshift\relax}}%
		\intgap
	}{% ветка без \limits: всё как раньше, без правок
		\IfValueT{#3}{_{#3}}%
		\IfValueT{#4}{^{#4}}%
	}%
}

\AtBeginDocument{%
	\let\intorig\int
	\let\ointorig\oint
	\renewcommand*{\int}{\opwithlimits{\intorig}}%
	\renewcommand*{\oint}{\opwithlimits{\ointorig}}%
}